\documentclass{article}
\usepackage[UTF8]{ctex}
\usepackage{amsmath,amssymb,amsfonts}
\usepackage{enumitem}
\usepackage{geometry}
\geometry{a4paper, margin=1in}
% === 错题本排版优化（不参与题目标识） ===
\setlength{\parindent}{0pt}
\setlength{\parskip}{0.35em plus 0.10em minus 0.08em}
\setlength{\emergencystretch}{3em}
\setlist{itemsep=0.25em, topsep=0.40em, parsep=0pt}
\allowdisplaybreaks[2]
% === 错题本排版优化结束 ===
\begin{document}
\section*{天津大学2026年数学分析真题}
\section*{60．天津大学2026年数学分析试卷}

\subsection{题目1}
1．（10 分）设 \(f\) 是 \(\mathbb{R}^{2}\) 上的连续函数，且 \(f\) 在原点沿任意方向的方向导数存在．问：\(f\) 在原点是否一定可微？若可微，给出证明，若不可微，给出反例．

\subsection{题目2}
2．（15 分）证明：\(f(x)=\sin \left(\frac{1}{x^{2}}\right)\) 在（0 ，1］上非一致连续，在 \([1,+\infty)\) 上一致连续．

\subsection{题目3}
3．（15 分）已知函数

\[
f(x)= \begin{cases}A, & 0 \leq x \leq \pi \\ -A, & -\pi \leq x<0\end{cases}
\]

其中 \(A \neq 0\) 为常数，求该函数在区间 \([-\pi, \pi]\) 的傅里叶级数，并证明 \(\sum_{n=1}^{\infty} \frac{1}{(2 n-1)^{2}}=\frac{\pi^{2}}{8}\) ．

\subsection{题目4}
4．（10 分）设函数 \(f(x)\) 在 \([0,+\infty)\) 的任何有限闭区间上可积，令 \(L(s)=\int_{0}^{+\infty} e^{-s x} f(x) \mathrm{d} x\) ，已知\(L\left(s_{0}\right)\) 收敛，证明：\(L(s)\) 在 \(s \geq s_{0}>0\) 时收敛。

\subsection{题目5}
5．（15 分）设 \(f(x) \in C[0,1]\) ，且 \(\int_{0}^{1} f(x) e^{n x} \mathrm{~d} x=0\) 对所有 \(n=0,1,2, \cdots\) 成立，证明：\(f(x) \equiv 0\) ．

\subsection{题目6}
6．（15 分）设 \(\mathbf{F}=\left(F^{1}, F^{2}, \cdots, F^{n}\right): \mathbb{R}^{n} \rightarrow \mathbb{R}^{n}\) 在原点可微，且 \(\mathbf{F}(0)=0\) 。已知 \(\sum_{i, j=1}^{n}\left(\frac{\partial F^{i}(0)}{\partial x_{j}}\right)^{2}<1\) ，设明：存在 \(\delta>0\) ，当 \(|\mathbf{x}| \leq \delta\) 时，有 \(|\mathbf{F}(\mathbf{x})| \leq \delta\) ．

\subsection{题目7}
7．（15 分）设函数 \(f, g: \mathbb{R} \rightarrow \mathbb{R}\) 连续可微，且 \(f(0)=0, f^{\prime}(0) \neq 0\) ，考虑方程 \(f(x)=t g(x)\)

\begin{itemize}
  \item[（1）] 证明：存在 \(\delta>0\) ，使得 \(x(t)\) 是方程在 \(t \in(-\delta, \delta)\) 上满足 \(x(0)=0\) 的唯一解．
  \item[（2）] 求 \(x(t)\) 在 \(t=0\) 处的一阶泰勒展开式．
\end{itemize}

\subsection{题目8}
8．（20 分）设 \(B\) 是 \(\mathbb{R}^{3}\) 中的单位开球，函数 \(u \in C^{2}(B) \cap C(\bar{B})\) 满足（i）\(u\) 是 \(B\) 上的非零函数，（ii） \(u \mid \partial B=0\) 。（iii）存在 \(\lambda \in \mathbb{R}\) ，使得在 \(B\) 上有 \(\Delta u=\lambda u\) 。

\begin{itemize}
  \item[（1）] 证明： \(\int_{B}|\nabla u|^{2} \mathrm{~d} x \mathrm{~d} y \mathrm{~d} z+\lambda \int_{B} u^{2} \mathrm{~d} x \mathrm{~d} y \mathrm{~d} z=0\) ．
  \item[（2）] 证明：\(\lambda<0\) ．
\end{itemize}

\subsection{题目9}
9．（20 分）设 \(\left\{b_{n}\right\}\) 是严格单调递增趋于 \(+\infty\) 的数列，且 \(\lim _{n \rightarrow \infty} \frac{a_{n}-a_{n-1}}{b_{n}-b_{n-1}}=A(A \in \mathbb{R})\) ．

\begin{itemize}
  \item[（1）] 证明： \(\lim _{n \rightarrow \infty} \frac{a_{n}}{b_{n}}=A\) ．
  \item[（2）] 计算 \(\lim _{n \rightarrow \infty} \frac{1^{k}+2^{k}+\cdots+n^{k}}{n^{k+1}}\) ，其中 \(k\) 为正整数．
\end{itemize}

\subsection{题目10}
10．（15 分）证明：自然数 \(e\) 是无理数．
\end{document}
